\documentclass[10pt,a4paper]{amsart}
\usepackage{amsmath,amssymb,amsthm,graphicx}
\usepackage{verbatim}
\newenvironment{reference}{\comment}{\endcomment}
%\newenvironment{reference}{}{}
\newenvironment{slogan}{\comment}{\endcomment}
\newenvironment{history}{\comment}{\endcomment}

% For commutative diagrams we use Xy-pic
\usepackage[all]{xy}

% We use 2cell for 2-commutative diagrams.
\xyoption{2cell}
% Theorem environments.
%
\theoremstyle{plain}
\newtheorem{theorem}[subsection]{Theorem}
\newtheorem{proposition}[subsection]{Proposition}
\newtheorem{lemma}[subsection]{Lemma}

\theoremstyle{definition}
\newtheorem{definition}[subsection]{Definition}
\newtheorem{example}[subsection]{Example}
\newtheorem{exercise}[subsection]{Exercise}
\newtheorem{situation}[subsection]{Situation}

\theoremstyle{remark}
\newtheorem{remark}[subsection]{Remark}
\newtheorem{remarks}[subsection]{Remarks}

\numberwithin{equation}{subsection}

% Macros
%
\def\lim{\mathop{\mathrm{lim}}\nolimits}
\def\colim{\mathop{\mathrm{colim}}\nolimits}
\def\Spec{\mathop{\mathrm{Spec}}}
\def\Hom{\mathop{\mathrm{Hom}}\nolimits}
\def\Ext{\mathop{\mathrm{Ext}}\nolimits}
\def\SheafHom{\mathop{\mathcal{H}\!\mathit{om}}\nolimits}
\def\SheafExt{\mathop{\mathcal{E}\!\mathit{xt}}\nolimits}
\def\Sch{\mathit{Sch}}
\def\Mor{\mathop{\mathrm{Mor}}\nolimits}
\def\Ob{\mathop{\mathrm{Ob}}\nolimits}
\def\Sh{\mathop{\mathit{Sh}}\nolimits}
\def\NL{\mathop{N\!L}\nolimits}
\def\CH{\mathop{\mathrm{CH}}\nolimits}
\def\proetale{{pro\text{-}\acute{e}tale}}
\def\etale{{\acute{e}tale}}
\def\QCoh{\mathit{QCoh}}
\def\Ker{\mathop{\mathrm{Ker}}}
\def\Im{\mathop{\mathrm{Im}}}
\def\Coker{\mathop{\mathrm{Coker}}}
\def\Coim{\mathop{\mathrm{Coim}}}

% Boxtimes
%
\DeclareMathSymbol{\boxtimes}{\mathbin}{AMSa}{"02}

%
% Macros for moduli stacks/spaces
%
\def\QCohstack{\mathcal{QC}\!\mathit{oh}}
\def\Cohstack{\mathcal{C}\!\mathit{oh}}
\def\Spacesstack{\mathcal{S}\!\mathit{paces}}
\def\Quotfunctor{\mathrm{Quot}}
\def\Hilbfunctor{\mathrm{Hilb}}
\def\Curvesstack{\mathcal{C}\!\mathit{urves}}
\def\Polarizedstack{\mathcal{P}\!\mathit{olarized}}
\def\Complexesstack{\mathcal{C}\!\mathit{omplexes}}
% \Pic is the operator that assigns to X its picard group, usage \Pic(X)
% \Picardstack_{X/B} denotes the Picard stack of X over B
% \Picardfunctor_{X/B} denotes the Picard functor of X over B
\def\Pic{\mathop{\mathrm{Pic}}\nolimits}
\def\Picardstack{\mathcal{P}\!\mathit{ic}}
\def\Picardfunctor{\mathrm{Pic}}
\def\Deformationcategory{\mathcal{D}\!\mathit{ef}}
\usepackage{cleveref}
\crefname{originallemma}{original lemma}{original lemmas}
\Crefname{originallemma}{Original lemma}{Original lemmas}
\makeatletter
\newlabel{categories-lemma-yoneda}{{001P}{source}{}{sourcepointer-1}{}}
\newlabel{categories-lemma-yoneda@cref}{{[originallemma][0][]001P}{[1][0][]source}{}{}{}}
\makeatother
\crefname{originallemma}{original lemma}{original lemmas}
\Crefname{originallemma}{Original lemma}{Original lemmas}
\makeatletter
\newlabel{lemma-sheaf-coequalizer-representable}{{00WS}{source}{}{sourcepointer-2}{}}
\newlabel{lemma-sheaf-coequalizer-representable@cref}{{[originallemma][0][]00WS}{[1][0][]source}{}{}{}}
\makeatother
\title{Existence of a left adjoint to the pullback of a morphism of topoi}
\author{The Stacks Project Authors}
\date{stacks-project:sites.tex:local-snapshot-20261012; GNU FDL 1.2 or later}
\begin{document}
\maketitle
\noindent\textit{Editorial note.} All delivered mathematical text is the author's own, copied from the pinned Stacks Project chapter file sites-cohomology.tex; no statement or proof step was rewritten, reordered or removed.\par\medskip
\noindent\textit{Editorial note.} The proof is the author's complete argument; every step the author writes out is inside the counted span and no closing sentence is a placeholder.\par\medskip
\noindent\textit{Editorial note.} Every labelled reference occurring inside the counted spans is carried in assembly.external\_references as a non-load-bearing navigation pointer (same-chapter ones at their definition site, cross-chapter ones at the invocation inside the counted span), so no load-bearing mathematics is replaced by a stub. No bibliography citation occurs inside the counted spans.\par\medskip
\noindent\textit{Editorial note.} The checked chapter list for same-name collisions is the set of stacks chapters staged in this repository: chow, cohomology, divisors, etale-cohomology, flat, more-algebra, more-morphisms, morphisms, spaces-more-morphisms, varieties together with this card's chapter sites-cohomology.\par\medskip
\noindent\textit{Editorial note.} The preamble is the chapter's own retained minimal preamble (sources/stacks-sites-cohomology/retained-preamble.tex), a verbatim excerpt of the native preamble: the Xy-pic setup and the native theorem declarations with the mathematical macros.\par\medskip
\noindent\textit{Allowed standard prerequisites (editorial).}\begin{itemize}
\item The Yoneda lemma. Categories, Lemma categories-lemma-yoneda. Shows that the representing object in the hypothesis is unique up to unique isomorphism and provides the identities used throughout.
\item Sheaves as coequalizers of coproducts of representables. Same chapter, Lemma lemma-sheaf-coequalizer-representable. Produces the coequalizer presentation of an arbitrary sheaf by coproducts of representable sheaves, using that the chosen family of objects covers.
\end{itemize}\medskip
\ifdefined\hypertarget\hypertarget{sourcepointer-1}{}\fi
\noindent\textit{Source reference.} Original lemma 001P: sources/stacks-sites/pinned-original-source.tex, source lines 5282--5282. Navigation pointer invoked at proof line 5282; the counted span contains the pointer only and the cited result is not reproduced inside it, so the pointer is registered as a non-load-bearing navigation stub in the established form for Stacks material. No bibliography citation is involved.\par\medskip
\ifdefined\hypertarget\hypertarget{sourcepointer-2}{}\fi
\noindent\textit{Source reference.} Original lemma 00WS: sources/stacks-sites/pinned-original-source.tex, source lines 5336--5336. Navigation pointer invoked at proof line 5336; the counted span contains the pointer only and the cited result is not reproduced inside it, so the pointer is registered as a non-load-bearing navigation stub in the established form for Stacks material. No bibliography citation is involved.\par\medskip
\begin{lemma}
\label{lemma-existence-lower-shriek}
Let $\mathcal{C}$, $\mathcal{D}$ be two sites.
Let $f : \Sh(\mathcal{D}) \to \Sh(\mathcal{C})$ be a morphism of topoi.
Let $E \subset \Ob(\mathcal{D})$ be a subset such that
\begin{enumerate}
\item for $V \in E$ there exists a sheaf $\mathcal{G}$
on $\mathcal{C}$ such that $f^{-1}\mathcal{F}(V) = 
\Mor_{\Sh(\mathcal{C})}(\mathcal{G}, \mathcal{F})$ functorially
for $\mathcal{F}$ in $\Sh(\mathcal{C})$,
\item every object of $\mathcal{D}$ has a covering by objects of $E$.
\end{enumerate}
Then $f^{-1}$ has a left adjoint $f_!$.
\end{lemma}

\begin{proof}
By the Yoneda lemma (Categories, Lemma \ref{categories-lemma-yoneda})
the sheaf $\mathcal{G}_V$ corresponding to $V \in E$
is defined up to unique isomorphism by the formula
$f^{-1}\mathcal{F}(V) = \Mor_{\Sh(\mathcal{C})}(\mathcal{G}_V, \mathcal{F})$.
Recall that
$f^{-1}\mathcal{F}(V) = \Mor_{\Sh(\mathcal{D})}(h_V^\#, f^{-1}\mathcal{F})$.
Denote $i_V : h_V^\# \to f^{-1}\mathcal{G}_V$ the map corresponding to
$\text{id}$ in $\Mor(\mathcal{G}_V, \mathcal{G}_V)$.
Functoriality in (1) implies that the bijection is given by
$$
\Mor_{\Sh(\mathcal{C})}(\mathcal{G}_V, \mathcal{F}) \to
\Mor_{\Sh(\mathcal{D})}(h_V^\#, f^{-1}\mathcal{F}),\quad
\varphi \mapsto f^{-1}\varphi \circ i_V
$$
For any $V_1, V_2 \in E$ there is a canonical map
$$
\Mor_{\Sh(\mathcal{D})}(h^\#_{V_2}, h^\#_{V_1})
\to
\Hom_{\Sh(\mathcal{C})}(\mathcal{G}_{V_2}, \mathcal{G}_{V_1}),\quad
\varphi \mapsto f_!(\varphi)
$$
which is characterized by
$f^{-1}(f_!(\varphi)) \circ i_{V_2} = i_{V_1} \circ \varphi$.
Note that $\varphi \mapsto f_!(\varphi)$ is
compatible with composition; this can be seen directly
from the characterization. Hence $h_V^\# \mapsto \mathcal{G}_V$
and $\varphi \mapsto f_!\varphi$ is a functor from
the full subcategory of $\Sh(\mathcal{D})$ whose objects are the $h_V^\#$.

\medskip\noindent
Let $J$ be a set and let $J \to E$, $j \mapsto V_j$ be a map.
Then we have a functorial bijection
$$
\Mor_{\Sh(\mathcal{C})}(\coprod \mathcal{G}_{V_j}, \mathcal{F})
\longrightarrow
\Mor_{\Sh(\mathcal{D})}(\coprod h_{V_j}^\#, f^{-1}\mathcal{F})
$$
using the product of the bijections above. Hence we can extend the
functor $f_!$ to the full subcategory of $\Sh(\mathcal{D})$ whose
objects are coproducts of $h_V^\#$ with $V \in E$.

\medskip\noindent
Given an arbitrary sheaf $\mathcal{H}$ on $\mathcal{D}$ we choose a
coequalizer diagram
$$
\xymatrix{
\mathcal{H}_1 \ar@<1ex>[r] \ar@<-1ex>[r] &
\mathcal{H}_0 \ar[r] &
\mathcal{H}
}
$$
where $\mathcal{H}_i = \coprod h_{V_{i, j}}^\#$
is a coproduct with $V_{i, j} \in E$.
This is possible by assumption (2), see
Lemma \ref{lemma-sheaf-coequalizer-representable}
(for those worried about set theoretical issues, note that
the construction given in
Lemma \ref{lemma-sheaf-coequalizer-representable} is canonical).
Define $f_!(\mathcal{H})$ to be the sheaf on $\mathcal{C}$
which makes
$$
\xymatrix{
f_!\mathcal{H}_1 \ar@<1ex>[r] \ar@<-1ex>[r] &
f_!\mathcal{H}_0 \ar[r] &
f_!\mathcal{H}
}
$$
a coequalizer diagram. Then
\begin{align*}
\Mor(f_!\mathcal{H}, \mathcal{F})
& =
\text{Equalizer}(
\xymatrix{
\Mor(f_!\mathcal{H}_0, \mathcal{F}) \ar@<1ex>[r] \ar@<-1ex>[r] &
\Mor(f_!\mathcal{H}_1, \mathcal{F})
}
) \\
& =
\text{Equalizer}(
\xymatrix{
\Mor(\mathcal{H}_0, f^{-1}\mathcal{F}) \ar@<1ex>[r] \ar@<-1ex>[r] &
\Mor(\mathcal{H}_1, f^{-1}\mathcal{F})
}
) \\
& =
\Hom(\mathcal{H}, f^{-1}\mathcal{F})
\end{align*}
Hence we see that we can extend $f_!$ to the whole category of sheaves
on $\mathcal{D}$.
\end{proof}

\end{document}
